\section{编者说明 · Editorial Note}

本书源于围绕 X Program \#161 开展的一系列数学与科研准备讨论，旨在系统理解递归自我改进（Recursive Self-Improvement）、统计学习理论以及金融研究环境中的样本外性能退化问题。

本书保留原始讨论的知识结构、数学推导、概念解释、研究动机、示例与可视化内容，并通过统一的 LaTeX 排版将其整理为适合长期阅读与反复查阅的个人学习材料。编辑过程以内容保真为原则，不以压缩篇幅或重新撰写教材为目标。

当前版本收录原始学习路线中已经完成的 Part I–IV。Part V 与 Part VI 属于后续研究计划，尚未形成完整的教学内容。

原始交互式可视化主要通过矢量图形及代表性参数状态呈现。无法直接恢复的原始资源、数据与引用均记录于配套的 Production Audit；编者重建的插图与原始图形予以明确区分。

本书中的理论推导、研究假设及实验设计主要用于科研准备与个人学习，不代表相关课题组的正式研究结论。内容保真审计亦不等同于对所有数学命题及经验性结论的独立学术验证。

% End of editorial note.
